% GENERATED FILE. Edit canonical lessons, then run npm run generate:books.
% canonical-source-hash: fnv1a64:314c969fb124247a
% canonical-lessons: TA-C370-varumai, TA-C370-celvam, TA-C370-errattalvu, TA-C370-velaiyinmai, TA-C370-curantu, TA-R370-review-363-370

\chapter{Rich and Poor}
\label{ch:ta-rich-and-poor}
\hlchaptermodality{370}

\begin{chapteropening}
\textbf{By the end of this chapter:} \emph{I can talk about poverty, wealth and inequality.}

Complete the last lesson of chapter 370: I can talk about poverty, wealth and inequality.
\end{chapteropening}

\section[\texorpdfstring{\ta{வறுமை} (vaṟumai)}{வறுமை (vaṟumai)}]{\ta{வறுமை} (vaṟumai) — poverty}

\label{lesson:TA-C370-varumai}



\begin{quote}
Before the new one: say the Tamil for a citizen, then the Tamil for to classify.
\end{quote}

\subsection*{You'll want to know: \ta{வறுமை}}
\textbf{\ta{வறுமை}} — \emph{vaṟumai} — \textquotedblleft{}poverty\textquotedblright{}.

\begin{quote}
\ta{பல} \ta{குடும்பங்கள்} \ta{இன்னும்} \ta{வறுமையில்} \ta{வாழ்கின்றன}.
\end{quote}

\begin{quote}
\emph{pala kuṭumpaṅkaḷ iṉṉum vaṟumaiyil vāḻkiṉṟaṉa.}
\end{quote}

\begin{quote}
Many families still live in poverty.
\end{quote}

\ta{வறுமையில்}, in poverty. \ta{ஏழை}, poor, is the word for the person.

\begin{cousinweb}
\textbf{\ta{வறுமை}} (\emph{vaṟumai}) is native Dravidian: DEDR 5513 lists it in the family of \emph{veṟu}, \textquotedblleft{}empty\textquotedblright{}, with \emph{veṟumai}, \textquotedblleft{}emptiness, poverty\textquotedblright{}: an empty hand.
\end{cousinweb}

\begin{grammarlens}[title={What you've built}]
One more word for giving reasons and opinions.
\end{grammarlens}

\subsection*{Your turn}
\begin{itemize}
\raggedright
  \item \emph{Say it:} \emph{vaṟumai}
  \item \emph{Say it:} \emph{vaṟumai}, once more
  \item \emph{Say it:} \emph{pala kuṭumpaṅkaḷ iṉṉum vaṟumaiyil vāḻkiṉṟaṉa.} — Many families still live in poverty.
  \item \emph{Recall:} say the Tamil for a citizen, then the Tamil for to classify, then say \emph{vaṟumai} again
\end{itemize}

\begin{tcolorbox}[breakable,colback=teal!4,colframe=teal!35!black,title={Before you move on}]
What does \ta{வறுமை} mean? (\textquotedblleft{}Poverty\textquotedblright{}.) Say it once more.
\end{tcolorbox}
\section[\texorpdfstring{\ta{செல்வம்} (celvam)}{செல்வம் (celvam)}]{\ta{செல்வம்} (celvam) — wealth; riches}

\label{lesson:TA-C370-celvam}



\begin{quote}
Before the new one: say the Tamil for to classify, then the Tamil for poverty.
\end{quote}

\subsection*{You'll want to know: \ta{செல்வம்}}
\textbf{\ta{செல்வம்}} — \emph{celvam} — \textquotedblleft{}wealth; riches\textquotedblright{}.

\begin{quote}
\ta{அவர்} \ta{தன்} \ta{செல்வத்தை} \ta{ஏழைகளுக்குக்} \ta{கொடுத்தார்}.
\end{quote}

\begin{quote}
\emph{avar taṉ celvattai ēḻaikaḷukkuk koṭuttār.}
\end{quote}

\begin{quote}
He gave his wealth to the poor.
\end{quote}

\ta{செல்வந்தர்}, a wealthy person.

\begin{cousinweb}
\textbf{\ta{செல்வம்}} (\emph{celvam}) is native Dravidian: DEDR 2786 lists Tamil \emph{celvam}, \textquotedblleft{}wealth, prosperity, happiness\textquotedblright{}, and \emph{celvaṉ}, \textquotedblleft{}a wealthy man\textquotedblright{}.
\end{cousinweb}

\begin{grammarlens}[title={What you've built}]
One more word for giving reasons and opinions.
\end{grammarlens}

\subsection*{Your turn}
\begin{itemize}
\raggedright
  \item \emph{Say it:} \emph{celvam}
  \item \emph{Say it:} \emph{celvam}, once more
  \item \emph{Say it:} \emph{avar taṉ celvattai ēḻaikaḷukkuk koṭuttār.} — He gave his wealth to the poor.
  \item \emph{Recall:} say the Tamil for to classify, then the Tamil for poverty, then say \emph{celvam} again
\end{itemize}

\begin{tcolorbox}[breakable,colback=teal!4,colframe=teal!35!black,title={Before you move on}]
What does \ta{செல்வம்} mean? (\textquotedblleft{}Wealth; riches\textquotedblright{}.) Say it once more.
\end{tcolorbox}
\section[\texorpdfstring{\ta{ஏற்றத்தாழ்வு} (ēṟṟattāḻvu)}{ஏற்றத்தாழ்வு (ēṟṟattāḻvu)}]{\ta{ஏற்றத்தாழ்வு} (ēṟṟattāḻvu) — inequality; ups and downs}

\label{lesson:TA-C370-errattalvu}



\begin{quote}
Before the new one: say the Tamil for poverty, then the Tamil for wealth.
\end{quote}

\subsection*{You'll want to know: \ta{ஏற்றத்தாழ்வு}}
\textbf{\ta{ஏற்றத்தாழ்வு}} — \emph{ēṟṟattāḻvu} — \textquotedblleft{}inequality; ups and downs\textquotedblright{}.

\begin{quote}
\ta{சமூகத்தில்} \ta{ஏற்றத்தாழ்வு} \ta{குறைய} \ta{வேண்டும்}.
\end{quote}

\begin{quote}
\emph{camūkattil ēṟṟattāḻvu kuṟaiya vēṇṭum.}
\end{quote}

\begin{quote}
Inequality in society must come down.
\end{quote}

Literally a rising and a falling (\ta{தாழ்வு}): highs and lows.

\begin{cousinweb}
\textbf{\ta{ஏற்றத்தாழ்வு}} (\emph{ēṟṟattāḻvu}) is native: \emph{ēṟṟam}, \textquotedblleft{}mounting, rising\textquotedblright{}, from \emph{ēṟu}, \textquotedblleft{}to rise\textquotedblright{} (DEDR 916), and \emph{tāḻvu}, \textquotedblleft{}lowness\textquotedblright{}, from \emph{tāḻ}, \textquotedblleft{}to be low\textquotedblright{} (DEDR 3178).
\end{cousinweb}

\begin{grammarlens}[title={What you've built}]
One more word for giving reasons and opinions.
\end{grammarlens}

\subsection*{Your turn}
\begin{itemize}
\raggedright
  \item \emph{Say it:} \emph{ēṟṟattāḻvu}
  \item \emph{Say it:} \emph{ēṟṟattāḻvu}, once more
  \item \emph{Say it:} \emph{camūkattil ēṟṟattāḻvu kuṟaiya vēṇṭum.} — Inequality in society must come down.
  \item \emph{Recall:} say the Tamil for poverty, then the Tamil for wealth, then say \emph{ēṟṟattāḻvu} again
\end{itemize}

\begin{tcolorbox}[breakable,colback=teal!4,colframe=teal!35!black,title={Before you move on}]
What does \ta{ஏற்றத்தாழ்வு} mean? (\textquotedblleft{}Inequality; ups and downs\textquotedblright{}.) Say it once more.
\end{tcolorbox}
\section[\texorpdfstring{\ta{வேலையின்மை} (vēlaiyiṉmai)}{வேலையின்மை (vēlaiyiṉmai)}]{\ta{வேலையின்மை} (vēlaiyiṉmai) — unemployment}

\label{lesson:TA-C370-velaiyinmai}



\begin{quote}
Before the new one: say the Tamil for wealth, then the Tamil for inequality.
\end{quote}

\subsection*{You'll want to know: \ta{வேலையின்மை}}
\textbf{\ta{வேலையின்மை}} — \emph{vēlaiyiṉmai} — \textquotedblleft{}unemployment\textquotedblright{}.

\begin{quote}
\ta{இந்த} \ta{ஊரில்} \ta{வேலையின்மை} \ta{ஒரு} \ta{பெரிய} \ta{பிரச்சினை}.
\end{quote}

\begin{quote}
\emph{inta ūril vēlaiyiṉmai oru periya piracciṉai.}
\end{quote}

\begin{quote}
Unemployment is a big problem in this town.
\end{quote}

\ta{வேலை}, work + \ta{இன்மை}, absence: -\ta{இன்மை} makes nouns of not having.

\begin{cousinweb}
\textbf{\ta{வேலையின்மை}} (\emph{vēlaiyiṉmai}) is built from \emph{vēlai}, \textquotedblleft{}work\textquotedblright{}, and native \emph{iṉmai}, \textquotedblleft{}non-existence, absence\textquotedblright{}, from \emph{il}, \textquotedblleft{}not to be\textquotedblright{} (DEDR 2559).
\end{cousinweb}

\begin{grammarlens}[title={What you've built}]
One more word for giving reasons and opinions.
\end{grammarlens}

\subsection*{Your turn}
\begin{itemize}
\raggedright
  \item \emph{Say it:} \emph{vēlaiyiṉmai}
  \item \emph{Say it:} \emph{vēlaiyiṉmai}, once more
  \item \emph{Say it:} \emph{inta ūril vēlaiyiṉmai oru periya piracciṉai.} — Unemployment is a big problem in this town.
  \item \emph{Recall:} say the Tamil for wealth, then the Tamil for inequality, then say \emph{vēlaiyiṉmai} again
\end{itemize}

\begin{tcolorbox}[breakable,colback=teal!4,colframe=teal!35!black,title={Before you move on}]
What does \ta{வேலையின்மை} mean? (\textquotedblleft{}Unemployment\textquotedblright{}.) Say it once more.
\end{tcolorbox}
\section[\texorpdfstring{\ta{சுரண்டு} (curaṇṭu)}{சுரண்டு (curaṇṭu)}]{\ta{சுரண்டு} (curaṇṭu) — to exploit; to scrape}

\label{lesson:TA-C370-curantu}



\begin{quote}
Before the new one: say the Tamil for inequality, then the Tamil for unemployment.
\end{quote}

\subsection*{You'll want to know: \ta{சுரண்டு}}
\textbf{\ta{சுரண்டு}} — \emph{curaṇṭu} — \textquotedblleft{}to exploit; to scrape\textquotedblright{}.

\begin{quote}
\ta{அந்த} \ta{முதலாளி} \ta{தொழிலாளர்களைச்} \ta{சுரண்டினார்}.
\end{quote}

\begin{quote}
\emph{anta mutalāḷi toḻilāḷarkaḷaic curaṇṭiṉār.}
\end{quote}

\begin{quote}
That employer exploited the workers.
\end{quote}

Past (I …): \ta{சுரண்டினேன்}, \emph{curaṇṭiṉēṉ}.

Who is exploited takes -\ta{ஐ}. Its first sense is to scrape, as a burnt pot is scraped clean.

\begin{cousinweb}
\textbf{\ta{சுரண்டு}} (\emph{curaṇṭu}) is native Dravidian: DEDR 1564 lists Tamil \emph{curaṇṭu}, \textquotedblleft{}to scrape\textquotedblright{}, in the family of \emph{cirai}, \textquotedblleft{}to shave, cut\textquotedblright{}; to exploit is to scrape someone bare.
\end{cousinweb}

\begin{grammarlens}[title={What you've built}]
One more word for giving reasons and opinions.
\end{grammarlens}

\subsection*{Your turn}
\begin{itemize}
\raggedright
  \item \emph{Say it:} \emph{curaṇṭu}
  \item \emph{Say it:} \emph{curaṇṭu}, once more
  \item \emph{Say it:} \emph{anta mutalāḷi toḻilāḷarkaḷaic curaṇṭiṉār.} — That employer exploited the workers.
  \item \emph{Recall:} say the Tamil for inequality, then the Tamil for unemployment, then say \emph{curaṇṭu} again
\end{itemize}

\begin{tcolorbox}[breakable,colback=teal!4,colframe=teal!35!black,title={Before you move on}]
What does \ta{சுரண்டு} mean? (\textquotedblleft{}To exploit; to scrape\textquotedblright{}.) Say it once more.
\end{tcolorbox}
\section[(dialogue)]{Looking back: from Numbers in an Argument to Rich and Poor}

\label{lesson:TA-R370-review-363-370}



\begin{quote}
Say the Tamil for statistics and a percentage.
\end{quote}

\subsection*{Your turn}
Say these aloud:

\begin{itemize}
\raggedright
  \item statistics, a percentage, an average, a ratio, to predict
  \item a professor, a lecture, a graduate, to study, admission
  \item an essay, a title, a preface, a summary, a quotation
  \item a field, psychology, geography, a formula, to define
  \item wisdom, ignorance, a scholar, a genius, to teach formally
  \item literacy, awareness, a scholarship, a seminar, a question
  \item society, the public, population, a citizen, to classify
  \item poverty, wealth, inequality, unemployment, to exploit
\end{itemize}

\begin{tcolorbox}[breakable,colback=teal!4,colframe=teal!35!black,title={Before you move on}]
Statistics? (\textbf{\ta{புள்ளிவிவரம்}}, \emph{puḷḷivivaram}.) A percentage? (\textbf{\ta{சதவீதம்}}, \emph{catavītam}.) An average? (\textbf{\ta{சராசரி}}, \emph{carācari}.) A ratio? (\textbf{\ta{விகிதம்}}, \emph{vikitam}.) To predict? (\textbf{\ta{கணி}}, \emph{kaṇi}.) A professor? (\textbf{\ta{பேராசிரியர்}}, \emph{pērāciriyar}.) A lecture? (\textbf{\ta{விரிவுரை}}, \emph{virivurai}.) A graduate? (\textbf{\ta{பட்டதாரி}}, \emph{paṭṭatāri}.) To study? (\textbf{\ta{பயில்}}, \emph{payil}.) Admission? (\textbf{\ta{சேர்க்கை}}, \emph{cērkkai}.) An essay? (\textbf{\ta{கட்டுரை}}, \emph{kaṭṭurai}.) A title? (\textbf{\ta{தலைப்பு}}, \emph{talaippu}.) A preface? (\textbf{\ta{முன்னுரை}}, \emph{muṉṉurai}.) A summary? (\textbf{\ta{சுருக்கம்}}, \emph{curukkam}.) A quotation? (\textbf{\ta{மேற்கோள்}}, \emph{mēṟkōḷ}.) A field? (\textbf{\ta{துறை}}, \emph{tuṟai}.) Psychology? (\textbf{\ta{உளவியல்}}, \emph{uḷaviyal}.) Geography? (\textbf{\ta{புவியியல்}}, \emph{puviyiyal}.) A formula? (\textbf{\ta{சூத்திரம்}}, \emph{cūttiram}.) To define? (\textbf{\ta{வரையறு}}, \emph{varaiyaṟu}.) Wisdom? (\textbf{\ta{ஞானம்}}, \emph{ñāṉam}.) Ignorance? (\textbf{\ta{அறியாமை}}, \emph{aṟiyāmai}.) A scholar? (\textbf{\ta{அறிஞர்}}, \emph{aṟiñar}.) A genius? (\textbf{\ta{மேதை}}, \emph{mētai}.) To teach formally? (\textbf{\ta{போதி}}, \emph{pōti}.) Literacy? (\textbf{\ta{கல்வியறிவு}}, \emph{kalviyaṟivu}.) Awareness? (\textbf{\ta{விழிப்புணர்வு}}, \emph{viḻippuṇarvu}.) A scholarship? (\textbf{\ta{உதவித்தொகை}}, \emph{utavittokai}.) A seminar? (\textbf{\ta{கருத்தரங்கு}}, \emph{karuttaraṅku}.) A question? (\textbf{\ta{வினா}}, \emph{viṉā}.) Society? (\textbf{\ta{சமூகம்}}, \emph{camūkam}.) The public? (\textbf{\ta{பொதுமக்கள்}}, \emph{potumakkaḷ}.) Population? (\textbf{\ta{மக்கள்தொகை}}, \emph{makkaḷtokai}.) A citizen? (\textbf{\ta{குடிமகன்}}, \emph{kuṭimakaṉ}.) To classify? (\textbf{\ta{வகைப்படுத்து}}, \emph{vakaippaṭuttu}.) Poverty? (\textbf{\ta{வறுமை}}, \emph{vaṟumai}.) Wealth? (\textbf{\ta{செல்வம்}}, \emph{celvam}.) Inequality? (\textbf{\ta{ஏற்றத்தாழ்வு}}, \emph{ēṟṟattāḻvu}.) Unemployment? (\textbf{\ta{வேலையின்மை}}, \emph{vēlaiyiṉmai}.) To exploit? (\textbf{\ta{சுரண்டு}}, \emph{curaṇṭu}.)
\end{tcolorbox}
